\chapter{联合 Monero 交易 (Joint Monero Transactions (TxTangle))}
\label{chapter:txtangle}

不同实体和场景的性质不可避免地产生了一些交易图启发式分析。特别是矿工、矿池（第 5.1 节，\cite{AnalysisOfLinkability}）、托管市场和交易所的行为在 Monero 基于环签名的协议中仍然有明显可分析的模式。

我们在此描述 TxTangle，类似于 Bitcoin 的 CoinJoin \cite{coinjoin-wiki}，这是一种混淆那些启发式的方法。\footnote{本章构成联合交易协议的提案。截至撰写本文时，尚未实现此类协议。之前的提案名为 MoJoin，由化名的合著者 Monero Research Lab (MRL) 研究员 Sarang Noether 创建，需要可信的经纪人才能运作。这种经纪人似乎与 Monero 项目对隐私和可替代性的基本承诺相矛盾，因此 MoJoin 未被进一步推进。}本质上，多笔交易被压缩成一笔交易，使每个参与者的行为模式融合在一起。

为了实现这种混淆，观察者利用联合交易中的信息来分组 input 和 output、将其与个别参与者关联、以及了解实际有多少参与者，必须是极其困难的。\footnote{由于在 Bitcoin 中金额是清晰可见的，通常可以根据金额总和对 CoinJoin 的 input 和 output 进行分组。\cite{coinjoin-sudoku}}此外，即使是参与者本身也不应知道参与者的数量，或者无法对其他参与者的 input 和 output 进行分组，除非他们控制了除一个参与者以外的所有参与者。\footnote{恶意污染联合交易是对这种方法的一种潜在攻击，最早针对 CoinJoin 被发现。\cite{coinjoin-pollution}}最后，应能在不依赖中心化机构的情况下构建联合交易 \cite{exa-blockchain-analysis}。幸运的是，在 Monero 中所有这些要求都可以实现。



\section{构建联合交易 (Building joint transactions)}
\label{sec:building-txtangle}

在普通交易中，input 和 output 通过金额平衡的证明绑定在一起。根据第 \ref{sec:commitments-and-fees} 节，伪 output 承诺之和等于 output 承诺之和（加上手续费承诺）。
\[\sum_j C'^a_{j} - (\sum_t C^b_{t} + f H) = 0\]

一个平凡的联合交易可以将多笔交易的所有内容放入一笔交易中。MLSAG 消息将签名所有子交易的数据，金额平衡也显然成立（0 + 0 + 0 = 0）。\footnote{由于 Bulletproofs 风格的范围证明实际上是聚合成一个的（第 \ref{sec:range_proofs} 节），即使在平凡情况下参与者也需要在一定程度上协作。}然而，input 和 output 的分组同样可以平凡地通过测试 input/output 子集的金额是否平衡来识别。\footnote{参与者可以尝试以花哨的方式分配手续费来混淆观察者，但面对暴力破解时会失败，因为手续费并不是很大（大约 32 位或更少）。}

我们可以通过计算每对参与者之间的共享秘密，然后将这些偏移量添加到他们的伪 output 承诺的掩码中来轻松解决这个问题（第 \ref{sec:ringct-introduction} 节）。在每一对中，一个参与者将共享秘密加到他的一个伪 output 承诺上，另一个参与者从{\em 他的}一个伪承诺中减去它。当求和时，秘密相互抵消，并且由于每对参与者都有一个共享秘密，金额平衡只有在所有承诺合并后才出现。\footnote{我们偏移的是伪 output 承诺而非 output 承诺，因为 output 承诺掩码是从接收者的地址构建的（第 \ref{sec:pedersen_monero} 节）。}

共享秘密可能在即时意义上隐藏 input/output 分组，但参与者必须以某种方式了解所有 input 和 output，最简单的方法是他们各自传达自己的 input/output 分组。这显然违反了初始前提，而且无论如何都意味着参与者知道参与者的数量。


\subsection{$n$ 路通信通道 ($n$-way communication channel)}
\label{subsec:n-way-channel}

TxTangle 的最大参与者数量是 output 或 input 的数量（取较小者）。我们通过每个真实参与者假装为他发送的每个 output 是不同的人来建模。这主要是为了与其他潜在参与者建立一个群组通信通道，而不暴露参与者的数量。

设想 $n$（$2 \leq n \leq 16$，但建议至少 3 人）\footnote{目前，一笔交易最多可以有 16 个 output。}个表面上互不相关的人在聊天室中以看似随机的间隔聚集，聊天室计划在时间 $t_0$ 开放、$t_1$ 关闭（一次只能有 16 个人在聊天室中，聊天室有条件如手续费优先级、每字节基础费用[用于围绕当前中位数和区块奖励简化共识]，以及可接受的交易类型范围——例如目前 Monero 创世初期的交易不能直接在 RingCT 交易中使用 \cite{pre-ringct-outputs-like-coinbase-research-issue-59}）。在 $t_1$ 时，所有模拟成员通过发布公钥来表示继续的意愿，然后通过在所有模拟成员之间构建共享秘密，聊天室被转换为 $n$ 路通信通道。\footnote{第 \ref{sec:m-of-n} 节中的 multisig 方法是一种方式，将 M-of-N 一路扩展到 1-of-N。}此共享秘密用于加密消息内容，而模拟成员使用 SAG 签名（第 \ref{SAG_section} 节）签名 input 相关消息，这样就永远不会清楚是谁发送了给定的消息，output 相关消息则用模拟成员公钥集合上的 bLSAG（第 \ref{blsag_note} 节）签名，使实际 output 与模拟成员脱钩。\footnote{每组独立的 TxTangle bLSAG 应使用相同的 key image，因为与给定 output 相关的所有内容都是链接在一起的。}


\subsection{构建联合交易的消息轮次 (Message rounds to construct a joint transaction)}
\label{subsec:message-rounds-txtangle}

通道设置完成后，TxTangle 交易可以在五轮通信中构建，下一轮只能在前一轮完成后才能开始，每轮都有一个定时通信间隔，消息应在其中随机发布。这些间隔旨在防止消息聚集暴露 input/output 分组。

\begin{enumerate}
    \item 每个模拟成员为每个预期 output 私下生成一个随机标量，并用 bLSAG 签名。这些标量的排序列表用于确定 output 索引（回忆第 \ref{sec:multi_out_transactions} 节；最小的标量获得索引 $t = 0$）。\footnote{output 索引选择应与其他交易构建实现匹配，以避免对不同软件进行指纹识别。我们使用这种有效随机的方法来与核心实现保持一致，后者也是随机化 output 的。}他们发布这些 bLSAG，以及签名预期 input 的交易版本号的 SAG。在这一轮之后，参与者可以根据 input 和 output 的数量计算交易权重，并准确估算所需的手续费。\footnote{如果根据比较 input 和 output 的数量与自己的 input/output 数量发现必须只有两个真实参与者，TxTangle 可以被放弃。建议每个参与者至少有两个 input 和两个 output，以防恶意行为者即使意识到只有两个参与者也不放弃 TxTangle。这个建议有待商榷，因为使用更多 input 和 output 在启发式上不是中性的。}\footnote{TxTangle 交易不应在 extra 字段中存储多余信息（例如不存储加密支付 ID，除非是只有 2 个 output 的 TxTangle，至少应有一个虚拟加密支付 ID）。}\footnote{手续费估算应基于标准化方法，使每个参与者计算出相同的结果。否则 output 可能会根据手续费计算方法被聚集。同样的手续费标准应在 TxTangle 之外实施，以促进 TxTangle 交易看起来与普通交易相同。}

    \item 每个模拟成员使用公钥列表与其他每个成员构建共享秘密来偏移他们的伪 output 承诺，并根据每对中较小的公钥决定谁加谁减。\footnote{由于点是压缩的（第 \ref{point_compression_section} 节），只需将密钥解释为 32 字节整数。按惯例，较小密钥的持有者加，较大密钥的持有者减。}每个模拟成员必须支付手续费估算的 $1/n$（使用整数除法）。拥有最小 output 索引的模拟成员负责支付除法后的余数（这将是一个极其微小的金额，但必须考虑以防止指纹识别 TxTangle 交易）。他们私下为每个 output 生成交易公钥（尚不发送给其他成员），并构建他们的 output 承诺、编码金额和用于聚合 Bulletproof 范围证明的 Part A 部分证明，用 bLSAG 签名所有这些（每个 bLSAG 消息一个承诺、一个编码金额和一个部分证明，key image 将这些与指定 output 索引的原始随机标量列表关联起来）。伪 output 承诺照常生成（第 \ref{sec:ringct-introduction} 节），然后用共享秘密偏移，并用 SAG 签名。bLSAG 和 SAG 发布后，假设参与者以相同方式估算了总手续费，他们现在可以验证金额整体上是否平衡。\footnote{我们不发布各自支付的单独手续费金额，以防某个参与者计算错误，这可能由于非标准手续费金额的集合而暴露 output 分组。如果金额不能正确平衡，TxTangle 交易可以被放弃。}

    \item 如果金额正确平衡，我们开始一个小的额外轮次来构建聚合 Bulletproof，以证明所有 output 金额在范围内。每个模拟成员使用上一轮的 Part A 部分证明和 output 承诺，并私下计算聚合挑战 A。他们用它来构建 Part B 部分证明，连同 bLSAG 一起发送到通道。

    \item 参与者开始填写将由 MLSAG 签名的消息（回忆第 \ref{full-signature} 节的脚注）。在通信间隔内按随机顺序发布的是两种消息。每个 input 的环成员偏移量和 key image 用 SAG 签名，并与正确的伪 output 承诺关联。每个 output 的一次性地址、交易公钥和 Part C 部分证明（基于 Part B 部分证明和聚合挑战 B 计算）用 bLSAG 签名（这些还可以包含一个随机的基交易公钥分量，正如我们将看到的，它可用于虚假 Janus 缓解）。

    \item 参与者使用所有部分证明完成聚合 Bulletproof，并私下应用对数内积技术将其压缩为最终证明以包含在交易数据中。一旦收集了所有将由 MLSAG 签名的信息，每个参与者完成他们的 input 的 MLSAG，并在通信间隔内随机将它们（每个带一个 SAG）发送到通道。任何参与者一旦获得所有片段就可以提交交易。
\end{enumerate}{}

\subsubsection*{交易公钥和 Janus 缓解 (Transaction public keys and the Janus mitigation)}

如果 TxTangle 的每个参与者都知道交易私钥 $r$（第 \ref{sec:one-time-addresses} 节），那么他们中的任何一个都可以根据已知地址列表测试其他人的一次性 output 地址。因此有必要将 TxTangle 交易构建为存在子地址接收者（第 \ref{sec:subaddresses} 节），包括为每个 output 使用不同的交易公钥。

为了匹配与子地址相关的 Janus 攻击缓解的可能实现——其中在 extra 字段中包含一个额外的「基」交易公钥 \cite{janus-mitigation-issue-62}，TxTangle 也应该有一个虚假的「基」密钥，由每个模拟成员生成的随机密钥之和组成。\footnote{密钥取消（第 \ref{subsec:drawbacks-naive-aggregation-cancellation} 节）应该不是问题，因为它只是一个虚假密钥，理想情况下应该在交易公钥列表中随机索引。}

许多向子地址汇款的 TxTangle 参与者可能至少有两个 output，其中一个将找零转回参与者。这意味着任何 TxTangle 参与者仍然可以通过将找零的交易公钥也作为子地址接收者的「基」密钥来启用 Janus 缓解。\footnote{如果向自己的子地址发送，则不需要 Janus 缓解。启用了 Janus 缓解的钱包应该识别出在 TxTangle 中花费的金额等于接收到子地址的金额，因此不会错误地通知用户存在问题。}子地址接收者可能会意识到交易是一个 TxTangle，并且「基」密钥可能对应于发送者的找零 output。\footnote{这是假设交易公钥与 output 是 1:1 的，这在目前看来是这种情况。如果交易公钥在 extra 字段中按随机或排序顺序是标准做法，那么 TxTangle 和非 TxTangle 交易对子地址接收者来说将是基本无法区分的。在某些特殊情况下，TxTangle 参与者无法包含「基」密钥（例如当他们的所有 output 都是发往子地址时），或者由于子地址接收者收到了大部分或全部 output 而明显不是 TxTangle。注意由于 TxTangle 交易通常比典型交易有更多的 output，此启发式可用于区分 TxTangle 和普通的子地址交易。}


\subsection{弱点 (Weaknesses)}
\label{subsec:weaknesses-txtangle}

恶意行为者有两种主要方式来破坏 TxTangle 的目的——即对潜在对手/分析师隐藏 input/output 分组。他们可以污染交易，使诚实参与者的子集更小（甚至不存在）\cite{coinjoin-pollution}。他们也可以导致 TxTangle 尝试失败，并利用相同参与者的后续尝试来估算 input/output 分组。

前者不容易缓解，特别是在非常去中心化的情况下，没有参与者有声誉。TxTangle 的一个可能应用是与协作的矿池一起使用，矿池可以隐藏其矿工属于哪个矿池。这种矿池会知道 input/output 分组，但由于目的是帮助其连接的矿工，保守信息对它们有利。此外，这种 TxTangle 不允许恶意行为者参与，假设矿池在某种程度上是诚实的。

后者可以通过仅尝试几次 TxTangle 后休息一下，并始终为新尝试重新生成交易的大多数随机元素来防御。这些元素包括交易公钥、伪承诺掩码、范围证明标量和 MLSAG 标量。特别是，每个 input 的环诱饵集应保持不变，以防止交叉比较暴露真正的 input。如果可能，不同的 TxTangle 尝试应使用不同的真正 input。由于这种弱点是不可避免的，它使下一节的概念更具吸引力。



\section{托管 TxTangle (Hosted TxTangle)}
\label{sec:hosted-txtangle}

真正去中心化的 TxTangle 有一些悬而未决的问题。定时轮次如何启动和执行？聊天室最初是如何创建的，让参与者找到彼此？最直接的方式是通过 TxTangle 主机，由它生成和管理那些聊天室。

这样的主机似乎违背了混淆参与者的目标，因为每个人必须连接并发送可用于关联 input/output 分组的消息（特别是如果主机参与并知道消息内容）。我们可以使用 I2P 等网络\footnote{隐形互联网项目 (\url{https://geti2p.net/en/})。}使主机收到的每条消息看起来像是来自一个独特个体。


\subsection{通过 I2P 与主机的基本通信及其他特性 (Basic communication with a host over I2P, and other features)}
\label{subsec:txtangle-host-communication}

使用 I2P，用户创建所谓的「隧道」，将重度加密的消息通过其他用户的客户端传送到目的地。据我们理解，这些隧道可以在被销毁和重新创建之前传输多条消息（例如隧道似乎有 10 分钟的定时器）。对于我们的用例，仔细控制新隧道的创建时间以及哪些消息可以从同一隧道传出是至关重要的。\footnote{在 I2P 中有「出站」和「入站」隧道（见 \url{https://geti2p.net/en/docs/how/tunnel-routing}）。通过入站隧道接收的所有内容看起来像是来自同一来源，即使来自多个来源，因此表面上看来 TxTangle 用户不需要为所有用例创建不同的隧道。然而，如果 TxTangle 主机将自己设置为其入站隧道的入口点，那么他就获得了对 TxTangle 参与者出站隧道的直接访问权。}

\begin{enumerate}
    \item {\em 申请 TxTangle}：在我们原始的 $n$ 路提案（第 \ref{subsec:n-way-channel} 节）中，参与者在计划关闭之前逐渐将他们的模拟成员添加到可用的 TxTangle 房间。然而，如果有足够多的用户同时尝试 TxTangle，当用户试图随机将所有预期 output 放入同一个 TxTangle 「房间」时，可能会有很高的失败率，然后房间过早满员，他们不得不退出。这会相当混乱。

    我们可以通过告诉主机我们有多少个 output（例如给他一个我们的模拟成员公钥列表），让他组装每个 TxTangle 的参与者，来实现有影响力的优化。由于我们仍然保留 bLSAG 和 SAG 消息协议，主机将无法识别最终交易中的 output 分组。他只知道参与者的数量以及每个人有多少 output。此外，在这种情况下观察者无法监控开放的 TxTangle 房间来推断参与者的信息，这是一个重要的隐私改进。注意主机污染 TxTangle 的能力与非主机设计没有显著差异，因此这一变化对该攻击向量是中性的。

    \item {\em 通信方法}：由于主机已经充当消息传输的中心，最简单的方式是由他管理 TxTangle 通信。在每一轮中，主机从模拟成员收集消息（仍然在通信间隔内随机进行），在轮次结束时有一个短暂的数据分发阶段，他将所有收集的数据发送给每个参与者，在下一轮之前有一个缓冲期以确保消息被接收并有时间处理。

    \item {\em 隧道和 input/output 分组}：一旦 TxTangle 启动，用户应将他们的模拟成员身份与实际创建的 output 分离。这意味着为 bLSAG 签名消息创建新隧道，每个此类隧道只能传输与特定 output 相关的消息（通过同一隧道传输多条此类消息是可以的，因为显然关于同一 output 的信息来自同一来源）。他们还应为与特定 input 相关的 SAG 签名消息创建新隧道。

    \item {\em 主机 MITM 攻击威胁}：主机可能通过假装是其他参与者来欺骗某个参与者，因为他控制着发送用于构建 bLSAG 和 SAG 的模拟成员列表。换言之，他发送给参与者 A 的列表可能包含参与者 A 的模拟成员，其余的都是他自己的。参与者 B 收到的消息在重新传输给 A 之前使用 A 的列表重新签名。由于 A 的列表签名的所有消息都属于 A，主机将直接了解 A 的 input/output 分组！

    我们可以通过修改交易公钥的制作方式来防止主机作为诚实参与者交互的 MITM。参与者照常互相发送他们的预期交易公钥（用 bLSAG），然后，很像第 \ref{sec:robust-key-aggregation} 节中的稳健密钥聚合，实际包含在交易数据中的密钥（用于制作 output 承诺掩码等）以模拟成员列表的哈希为前缀。换言之，$\mathcal{H}_n(T_{agg},\mathbb{S}_{mock},r_t G)*r_t G$ 是第 $t$ 个交易公钥。将模拟成员列表烘焙到交易本身中使得没有所有实际参与者之间的直接通信就很难完成 TxTangle。\footnote{如果实现了 Janus 缓解，此 MITM 防御应改为用虚假 Janus 基密钥实现。每个模拟成员提供一个随机密钥 $r_{mock} G$，然后实际基密钥为 $\sum_{mock} \mathcal{H}_n(T_{agg},\mathbb{S}_{mock},r_{mock} G)*r_{mock} G$。}
\end{enumerate}{}


\subsection{主机即服务 (Host as a service)}
\label{subsec:txtangle-host-service}

对于可持续性和持续改进，TxTangle 服务以营利方式运营很重要。\footnote{虽然部署的 TxTangle 服务可以营利，但代码本身可以是开源的。这对于审计与 TxTangle 服务交互的钱包软件很重要。}与其用基于账户的模型来损害用户身份，主机可以作为单独的 output 参与每个 TxTangle，并要求参与者为其提供资金。当访问主机的「eepsite」/服务申请 TxTangle 时，用户会收到当前托管费用的通知，该费用应按 output 支付。

参与者将负责支付手续费{\em 和}主机费用的份额。这次，最小模拟成员的公钥（不包括主机的密钥）将承担手续费和主机费用的余额。\footnote{此处必须使用模拟成员密钥，因为主机不支付手续费，且他的 output 索引未知。}由于主机没有 input，他没有伪 output 承诺来抵消其 output 的承诺掩码。相反，他照常与其他模拟成员创建共享秘密，然后将他真正的承诺掩码分成随机大小的块分配给每个其他模拟成员，并除以共享秘密。他发布这些标量的列表（根据公钥将它们与其他每个模拟成员对应），用他的模拟成员密钥签名，让参与者知道它来自主机。这个列表的出现标志着第 \ref{subsec:message-rounds-txtangle} 节中第 1 轮的开始（即设置轮「0」的结束）。模拟成员将他们的主机标量乘以适当的共享秘密，并加到他们的伪承诺掩码中。这样，即使是主机的 output 在最终交易中也不能被任何参与者识别，除非完全合谋对付他。

为了简化手续费的计算，主机可以在第 1 轮结束时分发交易中使用的总手续费，因为他会较早了解到交易权重。参与者可以验证该金额与预期金额相似，并支付其份额。

如果参与者合谋作弊不提供托管费用，主机可以在第 3 轮终止 TxTangle。如果通道中出现不该出现的或无效的消息，他也可以终止。

在第 5 轮结束时，主机完成交易并提交到网络进行验证，作为其服务的一部分。他在最终分发消息中包含交易哈希。



\section{使用可信经纪人 (Using a trusted dealer)}
\label{sec:dealer-txtangle}

去中心化 TxTangle 有一些缺点。它要求所有参与者在严格的时间表内积极通信（并首先找到彼此），而且实现起来具有挑战性。

添加一个负责从每个参与者收集交易信息并混淆 input/output 分组的中心化经纪人可以简化流程。代价是更高的信任要求，因为经纪人必须（至少）了解那些分组。\footnote{本节受 MoJoin 协议大纲的启发。}


\subsection{基于经纪人的流程 (Dealer-based procedure)}
\label{subsec:dealer-procedure-txtangle}

经纪人将宣传他可以管理 TxTangle，并收集潜在参与者的申请（包括预期 input [及其类型] 和 output 的数量）。如果他愿意，他可以用自己的 input/output 集参与。

在组装了接近 16 个 output 的分组后（应有两个或更多参与者，且没有参与者可以拥有除一个以外的所有 output 或 input），经纪人将开始五轮中的第一轮。在每一轮中，他从每个参与者收集信息，做出一些决定，并发出标志新轮次开始的消息。

\begin{enumerate}
    \item 首先，经纪人对每对参与者生成一个随机标量，并决定每对中谁应使用正版本或负版本。他使用 input 和 output 的数量和类型来估算所需的总手续费。他将每个参与者的标量求和，并私下向每个人发送其总和，加上他们预期支付的手续费份额，以及（随机选择的）他们的 output 索引。这些消息构成了 TxTangle 开始的信号。

    \item 每个参与者照常构建他们的子交易，为他们的 output 生成单独的交易公钥（根据需要进行 Janus 缓解），计算一次性 output 地址并编码 output 金额，创建平衡 output 承诺和手续费份额的伪 output 承诺，组装 MLSAG 签名中使用的环成员偏移量列表及相应的 key image，并将经纪人发送的标量（乘以 $G$）加到他们的一个伪 output 承诺上。他们为 output 创建 Part A 部分证明，并将所有这些信息发送给经纪人。经纪人验证 input 和 output 金额是否平衡，并将完整的 Part A 部分证明列表发送给每个参与者。

    \item 每个参与者计算聚合挑战 A，并生成 Part B 部分证明发送给经纪人。经纪人收集部分证明并分发给所有其他参与者。

    \item 每个参与者计算聚合挑战 B，并生成 Part C 部分证明发送给经纪人。经纪人收集这些证明，并应用对数内积技术将其压缩为最终证明。假设证明按预期验证，他生成一个随机的虚假 Janus「基」交易公钥，并将要由 MLSAG 签名的消息发送给每个参与者。

    \item 每个参与者完成他们的 MLSAG 并发送给经纪人。一旦他拥有了所有片段，他就可以完成交易的构建，并提交到区块链。他也可以将交易 ID 发送给每个参与者，以便他们确认交易已发布。
\end{enumerate}{}
